\documentclass[10pt,a4paper]{amsart}
\usepackage[margin=19mm]{geometry}
\usepackage{amsmath,amssymb,mathtools}
\usepackage[scr=rsfs,cal=euler]{mathalfa}
\usepackage{xfrac}
\usepackage[unicode,hidelinks,bookmarks=false]{hyperref}
\newcommand{\ie}{i.e.\ }
\newcommand{\cf}{cf.\ }
\newcommand{\resp}{respectively\ }
\newcommand{\Subsection}{\subsection}
\providecommand{\nobreakdash}{\nobreak\hbox{-}}
\setlength{\emergencystretch}{2em}
\multlinegap0pt
\usepackage{bm}
\usepackage{bbm}
\usepackage{dsfont}
\usepackage{xspace}
\usepackage{cancel}
\usepackage{mathtools}
\usepackage{amsthm}
\makeatletter
\@ifundefined{thm}{\newtheorem{thm}{Thm}}{}
\@ifundefined{lem}{\newtheorem{lem}[thm]{Lemma}}{}
\makeatother
\title[The characteristic polynomial of sunflowers]{The characteristic polynomial of sunflowers}
\author{Changjiang Bu, Lixiang Chen, Ge Lin}
\date{Source: arXiv:2506.17628v1 (CC BY 4.0)}
\begin{document}
\maketitle

\noindent\emph{Source and licence.} The characteristic polynomial of sunflowers. Version arXiv:2506.17628v1 (CC BY 4.0), distributed under the Creative Commons Attribution 4.0 International licence, as recorded in the arXiv licence metadata for this exact version and shown on the abstract page https://arxiv.org/abs/2506.17628v1. Full rights notice: licenses/arxiv-2506.17628-CC-BY-4.0.txt.\par\smallskip

\noindent\textit{Editorial scope.} Counted unit: the Lem whose source label is \texttt{yinli3.2} in the article (source0.tex lines 322--342, source label \texttt{yinli3.2}), delivered together with the article's own complete original proof of it (source0.tex lines 345--381, one \texttt{proof} environment of 37 lines). No mathematics is written, repaired or re-derived: statement and proof are the article's own text line for line. Required context retained uncounted: none. Every definition, display and label of those lines is retained unchanged. Omitted from this package and not redistributed: 1--321, 343--344, 382--810. Nothing inside a retained excerpt is omitted or altered: every retained body is delivered exactly as the article's lines are written, so no omission receipt of any kind accompanies this package. Citations: the retained excerpts carry no citation command at all, so no bibliography entry of the article is transcribed and no reference is redistributed. Reference adaptation: none was needed and none was made; every reference inside the delivered unit resolves inside it, so no unresolved reference or citation is produced. Printed slips kept literal: no printed word, symbol, hypothesis or number of the counted statement or of its proof was altered. Layout and class: the article's own class is \texttt{article} with packages including inputenc, subfig, graphicx, caption, color, amssymb, amsthm, amsmath, epic, setspace, float, natbib, multirow, hyperref; the wrapper is the lane's pinned \texttt{amsart} editorial wrapper at 10pt on A4, which restates the theorem-like environments the retained text needs and adds the article's own macro definitions that the retained text needs (none), with extra wrapper packages (bm, bbm, dsfont, xspace, cancel, mathtools). The delivered numbering is the unit's own. Assets: no retained span contains a figure environment, a graphics-inclusion command, a PDF-page inclusion, a table or a TikZ picture; no image file is copied into this package, the package declares no asset dependency and no image file is redistributed with it. Licence: version arXiv:2506.17628v1 of this article is distributed under the Creative Commons Attribution 4.0 International licence, as recorded in the arXiv licence metadata for this exact version and shown on the abstract page https://arxiv.org/abs/2506.17628v1. A full rights notice is delivered at licenses/arxiv-2506.17628-CC-BY-4.0.txt. Characters: every retained excerpt is pure ASCII, so the delivered engine, XeLaTeX with its default Latin Modern text font, reports no missing character.
\begin{lem}\label{yinli3.2}
For a sunflower $\mathcal{S}=\mathcal{S}(k,s,p)$,
let $S$ be the set of seed vertices, and let $P_{i}$ be the set of vertices on the $i$-th petal of $\mathcal{S}$ for all $i\in[p]$.
Let $d$ be a positive integer such that $k \mid d$ and $f \in \mathcal{F}_d(\mathcal{S})$.
%Let $f \in \mathcal{F}_d$.
Then

\begin{enumerate}
\renewcommand{\labelenumi}{(\alph{enumi})}
\item Let $\mathrm{deg}^+(v)$ and $\mathrm{deg}^-(v)$ be the out-degree and in-degree of the vertex $v$ in the multi-digraph $D_f$. Then $\mathrm{deg}^+(v) =\mathrm{deg}^-(v)$  for all $v \in V(D_f)$.

\item  For each $v \in S$ and each $i \in [p]$,  there exists a non-negative integer $q_{vi}$ such that $\{v\} \xrightarrow{q_{vi}}  P_{i}$ in $D_f$.

\item For each $i \in [p]$,  there exists a positive integer $m_i$ such that $P_{i} \xrightarrow{m_i} S \cup P_{i}$ in $D_f$.

\item When $s \geq 2$, then $S \xrightarrow{\frac{d}{k}} S$ in $D_f$ .
%\item There exists a positive integer $m$ such that $S \xrightarrow{m} S$ in $D_f$ if $s \geq 2$.


\end{enumerate}
\end{lem}

\begin{proof}

Note that the multi-digraph $D_f$ is Eulerian, it implies (a).

For each $i \in [p]$, let $e_i =S \cup P_i \in E(\mathcal{S})$.
For each $v \in S$, denote by $q_{vi}$ the number of $k$-tuples $v\alpha$ in $f$ such that the first entry is $v$ and $v\alpha=e_i$.
From the construction of $D_{f}$,  we have (b).


For each $i \in [p]$ and $u \in P_{i}$,
note that the $k$-tuples $u\alpha$ in $f$ with initial entry $u$ are exactly the hyperedge $e_i$ rooted at $u$, that is, $u\alpha=e_i$.
Let $m(v)$ denote the number of such $k$-tuples in $f$ with $v$ as their first entry for any $v \in V(D_{f})$.
For any $u \in P_{i}$,  we have $\mathrm{deg}^+(u)=(k-1)m(u)$
and $\mathrm{deg}^-(u)=\sum_{v\in P_{i}\setminus\{u\}}m(v)+\sum_{v\in S}q_{vi}$
(or $\mathrm{deg}^-(u)=\sum_{v\in S}q_{vi}$ if $s=k-1$).
Since $\mathrm{deg}^+(u)>0$, it follows that $m(u) >0$.
By (a),
we get
\begin{align}\label{eqpi}
km(u)=\sum_{v\in P_{i}}m(v)+\sum_{v\in S}q_{vi},
\end{align}
where the right-hand side is independent of $u$. It implies that $m(u)$ is constant for all $u \in P_{i}$.
For each $i \in [p]$ and all $u \in P_{i}$, we set $m(u)=m_{i}$.
From the construction of $D_{f}$, we have (c).



For each $v \in S$, observe that $v$ appears in every rooted hyperedge in $f$.
Note that $m(v)$ counts the number of such hyperedges rooted at $v$,
and thus $d-m(v)$ counts the number of such hyperedges not rooted at $v$.
Then we have $\mathrm{deg}^+(v)=(k-1)m(v)$ and $\mathrm{deg}^-(v)=d-m(v)$.
%Because $v$ is contained in all $k$-tuples in $f$ that do not have it as the first entry,
%it follows that $\mathrm{deg}^-(v)=d-m_{S}(v)$.
By (a),
we get $m(v)=\frac{d}{k}$ for each $v \in S$.
When $s \geq 2$, for any two distinct vertices $v_1,v_2 \in S$, the multiplicity of arc in $D_f$ from $v_1$ to $v_2$ is $m(v_1)=\frac{d}{k}$, it implies (d).
\end{proof}

\end{document}
